% Reparto del conjunto de evaluacion (IT-27, ampliado en IT-100 y en IT-86).
% Cifras y ejemplos literales de eval/preguntas_evaluacion.json, comprobados
% con `py scripts/check_evalset.py`: 66 preguntas (56 de dominio y 10 ajenas),
% 11/11 titulaciones alcanzadas de las 12 del corpus.
%
% NO ajustar estas cifras a mano: si no coinciden, se ejecuta el verificador y
% se copia su salida. La tabla decia 50 preguntas y 2 sin guia mientras el
% fichero versionado tenia 66 y 8, porque no se refresco al ampliar el conjunto.
\begin{table}[htbp]
  \centering
  \small
  \caption{Reparto del conjunto de evaluación por tipo de pregunta. Fuente: elaboración propia.}
  \label{tab:conjunto_evaluacion}
  \begin{tabularx}{\textwidth}{@{}l X r@{}}
    \toprule
    \textbf{Tipo} & \textbf{Ejemplo}                                                     & \textbf{n} \\
    \midrule
    Temario       & ¿Qué temario tiene Estructuras de datos?                              & 20         \\
    Listado       & ¿Qué materias son obligatorias en el Grado en Ingeniería Mecánica?    & 14         \\
    Salidas       & ¿Qué trabajos hay para alguien que estudie Ingeniería Mecánica?       & 8          \\
    \textcolor{borrador2}{Sin guía} & \textcolor{borrador2}{¿Qué se estudia en Aprendizaje profundo del grado de IA y Ciberseguridad?} & \textcolor{borrador2}{8} \\
    Metadatos     & ¿Cuántos ECTS tiene Diseño de software?                               & 6          \\
    \midrule
    \textcolor{borrador2}{\textbf{De dominio}} & & \textcolor{borrador2}{\textbf{56}} \\
    \midrule
    \textcolor{borrador2}{Fuera de dominio} & \textcolor{borrador2}{¿Puedo estudiar Medicina en la Escuela Politécnica Superior de Jaén?} & \textcolor{borrador2}{10} \\
    \midrule
    \textbf{Total} &                                                                     & \textcolor{borrador2}{\textbf{66}} \\
    \bottomrule
  \end{tabularx}
  \begin{minipage}{0.95\textwidth}
    \vspace{4pt}
    \tiny
    \textbf{Notas:} los ejemplos son literales del conjunto. Cada tipo apunta a un origen
    distinto de la colección: temario y metadatos, a la guía docente de la asignatura;
    salidas, al bloque de salidas de la titulación; sin guía, al fragmento informativo de una
    asignatura cuya guía no está publicada; y listado, a la unidad de plan de estudios, que es
    la que contiene el listado ya agregado. \textcolor{borrador2}{Las diez preguntas
      \textbf{fuera de dominio} no apuntan a ningún fragmento y se cuentan aparte: su lista de
      relevantes está vacía, de modo que no entran en el Recall@K ni en el MRR ---aportarían un
      cero fijo a los dos--- y su criterio de acierto es el contrario, no recuperar nada.}
  \end{minipage}
\end{table}
